\documentclass[12pt]{article}
\usepackage[utf8]{inputenc}
\usepackage{graphicx}
\usepackage{hyperref}
\usepackage{listings}
\usepackage{xcolor}
\usepackage{geometry}
\usepackage{enumitem}
\usepackage{float}

\geometry{a4paper, margin=1in}

\hypersetup{
    colorlinks=true,
    linkcolor=blue,
    filecolor=magenta,      
    urlcolor=cyan,
    pdftitle={Phase 3 Report: Enhanced Indexing, Fault Tolerance, and Monitoring},
    pdfauthor={Team},
}

\definecolor{codegreen}{rgb}{0,0.6,0}
\definecolor{codegray}{rgb}{0.5,0.5,0.5}
\definecolor{codepurple}{rgb}{0.58,0,0.82}
\definecolor{backcolour}{rgb}{0.95,0.95,0.92}
\definecolor{mydarkgreen}{RGB}{0,125,0}

\lstdefinestyle{mystyle}{
    backgroundcolor=\color{backcolour},   
    commentstyle=\color{codegreen},
    keywordstyle=\color{magenta},
    numberstyle=\tiny\color{codegray},
    stringstyle=\color{codepurple},
    basicstyle=\ttfamily\footnotesize,
    breakatwhitespace=false,         
    breaklines=true,                 
    captionpos=b,                    
    keepspaces=true,                 
    numbers=left,                    
    numbersep=5pt,                  
    showspaces=false,                
    showstringspaces=false,
    showtabs=false,                  
    tabsize=2
}

\lstset{style=mystyle}

\title{Distributed Computing - CSE354 \\ Spring 2025\\Project - Phase 3 Report\\Enhanced Indexing, Fault Tolerance, and Monitoring}

\author{Presented to : Dr. Ayman Bahaa, Dr. Hossam Abdelrahman,\\ Eng. Mostafa Ashraf, Eng. A'laa Hamdy}

\date{}

\begin{document}

\maketitle

\begin{center}
    \includegraphics[width=0.7\textwidth]{logoo (2) (1).PNG} % Adjust the width as needed
\end{center}

\begin{table}[htp]
\begin{center}
\caption{Team Members}
\label{tab:team10}
\vspace{10pt} % Adjust the value to increase or decrease the space
\begin{tabular}{|c|c|}
\hline
\textbf{ID} & \textbf{Team Member} \\
\hline
21P0180 & Patrick Ramez Eskander Bolous\\
\hline
21P0173 & Monica Hany Makram Derias \\
\hline
2100561 & Ahmed Tarek Mahmoud\\
\hline 
21P0082 & Youssef Hany Emil Ayad\\
\hline
21P0258 & Youssef Adel Albert\\
\hline
21P0206 & Mark Saleh Sobhi \\
\hline
\end{tabular}
\end{center}
\end{table}

\newpage
\tableofcontents
\newpage

\section{Introduction}
This report details the implementation of Phase 3 of our Distributed Web Crawling and Indexing System. In this phase, we have enhanced the system with improved indexing capabilities, fault tolerance mechanisms, and comprehensive monitoring. The focus has been on improving the robustness and reliability of the system while adding more advanced search functionality.

Building upon the foundation established in Phase 2, we have integrated advanced indexing using the Whoosh library, implemented fault tolerance for Crawler Nodes through heartbeat monitoring and task re-queueing, developed a monitoring dashboard for system metrics, and added data persistence through cloud storage. We have also created a search interface that allows users to submit complex queries with boolean operators and phrase searching.

This report provides a comprehensive overview of all the enhancements implemented in Phase 3, including implementation details, architectural decisions, and testing results. We will document how our system meets the Phase 3 requirements outlined in the project specifications and discuss the challenges faced during implementation along with their solutions. 

\section{System Architecture}

\subsection{Enhanced System Architecture Overview}
The architecture of our distributed web crawling and indexing system has been enhanced in Phase 3 to include fault tolerance, advanced indexing, and monitoring capabilities. The system maintains its master-worker architecture with the following components:

\begin{itemize}
    \item \textbf{Master Node}: Coordinates the crawling process, distributes tasks, monitors Crawler Node health through heartbeats, detects failures, and re-queues tasks from failed nodes.
    \item \textbf{Crawler Nodes}: Fetch web pages, extract content and URLs, send regular heartbeat signals to the Master Node, and handle occasional simulated failures to test fault tolerance.
    \item \textbf{Enhanced Indexer Node}: Uses Whoosh for robust indexing with advanced text processing features including stemming, keyword extraction, and summarization.
    \item \textbf{Monitoring Dashboard}: A web-based interface that displays real-time system metrics, crawler status, and performance data.
    \item \textbf{Search Interface}: A web interface allowing users to perform advanced searches using boolean operators and phrase queries.
    \item \textbf{Cloud Storage}: AWS S3 integration for durable storage of both raw HTML and processed text data.
\end{itemize}

Figure 1 illustrates the enhanced architecture with the new components and interactions added in Phase 3.

\subsection{Fault Tolerance Mechanisms}
The system now implements multiple fault tolerance mechanisms to handle Crawler Node failures:

\begin{enumerate}
    \item \textbf{Heartbeat Monitoring}: Each Crawler Node sends regular heartbeat signals to the Master Node (every 2-5 seconds) to indicate it is still active and healthy.
    \item \textbf{Task Timeout Detection}: The Master Node tracks the start time of each task and detects if a task exceeds a specified timeout threshold (default: 60 seconds).
    \item \textbf{Task Re-queueing}: When a node fails (either through missed heartbeats or task timeouts), the Master Node re-queues its assigned task for processing by other active nodes.
    \item \textbf{Node Status Tracking}: The system maintains and updates the status of each Crawler Node (active or failed) in the system metrics.
\end{enumerate}

\subsection{Monitoring System}
A comprehensive monitoring system has been implemented with the following components:

\begin{itemize}
    \item \textbf{Enhanced Logging}: All system components (Master, Crawler, and Indexer Nodes) now provide detailed logging to track operations, errors, and performance metrics.
    \item \textbf{Centralized Metrics Collection}: The Master Node collects and maintains system-wide metrics including URLs crawled, indexed, failed, crawler status, and error rates.
    \item \textbf{Monitoring Dashboard}: A Flask-based web application that provides real-time visualization of system metrics, crawler node status, and performance data.
    \item \textbf{Monitoring Data Storage}: System metrics are regularly written to a JSON file (\texttt{monitoring\_data.json}) which is read by the dashboard application.
\end{itemize}

\subsection{Data Flow and Component Interaction}
The enhanced data flow in Phase 3 includes the following interactions:

\begin{enumerate}
    \item \textbf{Master Node Operations}:
    \begin{itemize}
        \item Assigns tasks to active Crawler Nodes
        \item Monitors heartbeats from Crawler Nodes
        \item Detects task timeouts and node failures
        \item Re-queues tasks from failed nodes
        \item Updates monitoring metrics and saves them to \texttt{monitoring\_data.json}
    \end{itemize}
    
    \item \textbf{Crawler Node Operations}:
    \begin{itemize}
        \item Sends regular heartbeat signals to the Master Node
        \item Fetches and processes web pages
        \item Extracts URLs and sends them to the Master Node
        \item Sends content to the Indexer Node
        \item Sends status updates and error reports to the Master Node
        \item Simulates occasional failures (5\% chance) to test fault tolerance
    \end{itemize}
    
    \item \textbf{Indexer Node Operations}:
    \begin{itemize}
        \item Processes received content using advanced NLP techniques
        \item Stores both raw HTML and processed text in cloud storage
        \item Builds and maintains a Whoosh-based search index
        \item Sends indexing status updates to the Master Node
    \end{itemize}
    
    \item \textbf{Monitoring Dashboard}:
    \begin{itemize}
        \item Reads system metrics from \texttt{monitoring\_data.json}
        \item Reads log files to determine component status
        \item Updates its display every 2 seconds with current system state
        \item Provides API endpoints for accessing system metrics
    \end{itemize}
    
    \item \textbf{Search Interface}:
    \begin{itemize}
        \item Provides a web interface for users to submit search queries
        \item Supports advanced search features including boolean operators and phrase search
        \item Displays search results with relevance scores and metadata
    \end{itemize}
\end{enumerate} 

\section{Implementation Details}

\subsection{Enhanced Indexer Implementation}
In Phase 3, we have significantly improved the indexing capabilities by implementing a robust indexing system using the Whoosh library. The enhanced indexer provides advanced text processing, efficient search functionality, and integration with cloud storage for data persistence.

\subsubsection{Whoosh-based Indexing}
We implemented a Whoosh-based indexing system with a rich schema that supports advanced search features:

\begin{lstlisting}[language=Python, caption=Whoosh Schema Implementation]
# Define the schema for our index
self.schema = Schema(
    url=ID(stored=True, unique=True),
    title=TEXT(stored=True, analyzer=StemmingAnalyzer()),
    content=TEXT(stored=True, analyzer=StemmingAnalyzer()),
    keywords=TEXT(stored=True),
    summary=STORED,
    last_updated=STORED
)
\end{lstlisting}

The schema includes:
\begin{itemize}
    \item URL field as a unique identifier
    \item Title and content fields with stemming analysis for improved search
    \item Keywords field for storing extracted key terms
    \item Summary field for storing a brief summary of the content
    \item Last updated timestamp for tracking freshness
\end{itemize}

\subsubsection{Advanced Text Processing}
The indexer implements advanced text processing techniques using NLTK for improved content analysis:

\begin{lstlisting}[language=Python, caption=Advanced Text Processing]
def process_text(self, text):
    """Process text with NLP techniques"""
    try:
        # Tokenize into sentences and words
        sentences = sent_tokenize(text.lower())
        
        # Process each sentence
        processed_words = []
        for sentence in sentences:
            words = word_tokenize(sentence)
            # Filter and lemmatize words
            words = [self.lemmatizer.lemmatize(word) for word in words 
                    if word.isalnum() and word not in self.stop_words]
            processed_words.extend(words)
            
        return processed_words
    except Exception as e:
        # Fallback to basic word splitting if NLTK fails
        logging.warning(f"NLTK processing failed, using basic tokenization: {str(e)}")
        words = text.lower().split()
        return [word.strip('.,!?()[]{}:;"\'') for word in words 
                if word.strip('.,!?()[]{}:;"\'').isalnum() 
                and len(word) > 3 
                and word not in self.stop_words]
\end{lstlisting}

Key text processing features include:
\begin{itemize}
    \item Sentence and word tokenization
    \item Stopword removal to filter out common words
    \item Lemmatization to reduce words to their root form
    \item Fallback mechanism for handling NLTK processing errors
\end{itemize}

\subsubsection{Keyword Extraction and Summarization}
The indexer implements keyword extraction and content summarization to enhance search results:

\begin{lstlisting}[language=Python, caption=Keyword Extraction and Summarization]
def extract_keywords(self, processed_words, top_n=10):
    """Extract key terms based on frequency"""
    from collections import Counter
    word_freq = Counter(processed_words)
    return [word for word, _ in word_freq.most_common(top_n)]

def generate_summary(self, title, text, max_sentences=3):
    """Generate a brief summary using key sentences"""
    sentences = sent_tokenize(text)
    if not sentences:
        return ""
        
    # Simple extractive summarization
    if len(sentences) <= max_sentences:
        return " ".join(sentences)
        
    # Use first sentence (usually most important) and last few based on max_sentences
    summary = [sentences[0]]
    if max_sentences > 1:
        summary.extend(sentences[-(max_sentences-1):])
    return " ".join(summary)
\end{lstlisting}

\subsubsection{Enhanced Search Functionality}
The indexer provides advanced search capabilities including boolean operators, phrase search, and field-specific search:

\begin{lstlisting}[language=Python, caption=Enhanced Search Implementation]
def search(self, query_string, search_fields=None, page=1, pagelen=10):
    """
    Enhanced search functionality supporting:
    - Multi-field search
    - Boolean operators (AND, OR, NOT)
    - Phrase search (using quotes)
    - Field-specific search (field:term)
    """
    try:
        search_fields = search_fields or ["title", "content", "keywords"]
        
        with self.ix.searcher(weighting=scoring.BM25F) as searcher:
            parser = MultifieldParser(search_fields, self.ix.schema)
            query = parser.parse(query_string)
            
            results = searcher.search_page(query, page, pagelen=pagelen)
            
            return {
                "total_results": len(results),
                "current_page": page,
                "results": [
                    {
                        "url": result["url"],
                        "title": result["title"],
                        "summary": result["summary"],
                        "keywords": result["keywords"],
                        "score": result.score,
                        "last_updated": result["last_updated"]
                    }
                    for result in results
                ]
            }
            
    except Exception as e:
        return {
            "status": "error",
            "error": str(e)
        }
\end{lstlisting}

The search functionality includes:
\begin{itemize}
    \item Multi-field search across title, content, and keywords
    \item Boolean operators (AND, OR, NOT) for complex queries
    \item Phrase search using quotes for exact matching
    \item Field-specific search using field:term syntax
    \item Pagination support for handling large result sets
    \item BM25F scoring algorithm for relevant result ranking
\end{itemize} 

\subsection{Fault Tolerance Implementation}
In Phase 3, we have implemented comprehensive fault tolerance mechanisms to handle Crawler Node failures and ensure continuous system operation. The fault tolerance system is primarily managed by the Master Node and includes heartbeat monitoring, task timeout detection, and task re-queueing.

\subsubsection{Heartbeat Monitoring System}
Each Crawler Node sends regular heartbeat signals to the Master Node to indicate it is still active:

\begin{lstlisting}[language=Python, caption=Crawler Heartbeat Implementation]
# Heartbeat mechanism in crawlerNode.py
def send_heartbeat():
    """Send periodic heartbeats to the master node"""
    while not shutdown_flag:
        try:
            uptime = time.time() - stats["urls_processed"]
            heartbeat_msg = {
                "type": "heartbeat",
                "rank": rank,
                "uptime": uptime,
                "urls_processed": stats["urls_processed"],
                "errors": stats["errors"]
            }
            comm.send(f"Heartbeat from crawler {rank}: Active", dest=0, tag=99)
            logging.debug(f"Sent heartbeat to master node")
        except Exception as e:
            logging.error(f"Error sending heartbeat: {e}")
        
        # Sleep for a random time between 2 and 5 seconds to avoid synchronized heartbeats
        time.sleep(random.uniform(2, 5))

# Start heartbeat thread
heartbeat_thread = threading.Thread(target=send_heartbeat, daemon=True)
heartbeat_thread.start()
\end{lstlisting}

The Master Node monitors these heartbeats and tracks the last heartbeat time for each Crawler Node:

\begin{lstlisting}[language=Python, caption=Master Heartbeat Monitoring]
# In masterNode.py
# Process heartbeat message
if tag == 99:
    source = status.Get_source()
    heartbeat_timestamps[source] = datetime.now()
    
    if "Heartbeat" in message:
        # Update crawler status if it was previously marked as failed
        if system_metrics["crawler_status"].get(str(source)) == "failed":
            logging.info(f"Crawler {source} has recovered.")
            system_metrics["crawler_status"][str(source)] = "active"
    else:
        # Process other status update messages
        logging.info(f"Status update from node {source}: {message}")
\end{lstlisting}

\subsubsection{Task Timeout Detection}
The Master Node tracks the start time of each task and checks for timeouts during each iteration of its main loop:

\begin{lstlisting}[language=Python, caption=Task Timeout Detection]
# Check for task timeouts
current_time = datetime.now()
for node, assignment in list(crawler_assignments.items()):
    if assignment is not None:
        # Check if task has exceeded the timeout threshold
        if task_start_times[node] and (current_time - task_start_times[node]).total_seconds() > task_timeout_seconds:
            logging.warning(f"Task timeout for crawler {node} on URL: {assignment}")
            # Re-queue the URL
            task_queue.send_message(assignment, {"type": "retry", "failed_node": str(node)})
            system_metrics["urls_failed"] += 1
            system_metrics["error_count"] += 1
            
            # Update stats for the assignment
            crawler_assignments[node] = None
            task_start_times[node] = None
            
            # Log the timeout for monitoring
            system_metrics["task_failures"].append({
                "url": assignment,
                "node": str(node),
                "time": current_time.isoformat(),
                "reason": "task_timeout"
            })
\end{lstlisting}

\subsubsection{Heartbeat Timeout Detection}
The Master Node also checks for missed heartbeats to detect failed Crawler Nodes:

\begin{lstlisting}[language=Python, caption=Heartbeat Timeout Detection]
# Check for heartbeat timeouts
for node in active_crawler_nodes:
    if node in heartbeat_timestamps:
        # Check if heartbeat has exceeded the timeout threshold
        if (current_time - heartbeat_timestamps[node]).total_seconds() > heartbeat_timeout_seconds:
            logging.warning(f"Heartbeat timeout for crawler {node}. Marking as failed.")
            # Mark node as failed in system metrics
            system_metrics["crawler_status"][str(node)] = "failed"
            
            # Re-queue any assigned task
            if crawler_assignments[node] is not None:
                assignment = crawler_assignments[node]
                task_queue.send_message(assignment, {"type": "retry", "failed_node": str(node)})
                system_metrics["urls_failed"] += 1
                system_metrics["error_count"] += 1
                
                # Log the failure for monitoring
                system_metrics["task_failures"].append({
                    "url": assignment,
                    "node": str(node),
                    "time": current_time.isoformat(),
                    "reason": "heartbeat_timeout"
                })
                
                # Clear the assignment
                crawler_assignments[node] = None
                task_start_times[node] = None
\end{lstlisting}

\subsubsection{Task Re-queueing}
When a failure is detected, the Master Node re-queues the affected task for processing by other active nodes:

\begin{lstlisting}[language=Python, caption=Task Re-queueing Implementation]
# Re-queue a task after failure
task_queue.send_message(assignment, {"type": "retry", "failed_node": str(node)})

# Function to assign tasks from the queue to available crawlers
def assign_tasks_from_queue():
    """Assign tasks from the queue to available crawlers"""
    for node in active_crawler_nodes:
        if crawler_assignments[node] is None and system_metrics["crawler_status"][str(node)] == "active":
            # Retrieve a message from the queue
            messages = task_queue.receive_messages(max_messages=1, wait_time=0)
            if messages:
                message = messages[0]
                # Extract the URL from the message body
                url_to_crawl = message.body if hasattr(message, 'body') else message.get('MessageBody', '')
                
                # Delete the message from the queue
                task_queue.delete_message(message)
                
                # Assign URL to crawler
                crawler_assignments[node] = url_to_crawl
                task_start_times[node] = datetime.now()
                
                # Send the URL to the crawler for processing
                comm.send(url_to_crawl, dest=node, tag=0)
                
                # Log the assignment for monitoring
                system_metrics["task_assignments"].append({
                    "url": url_to_crawl,
                    "node": str(node),
                    "time": datetime.now().isoformat(),
                    "type": "new_assignment"
                })
                
                logging.info(f"Assigned {url_to_crawl} to crawler {node}")
\end{lstlisting}

\subsubsection{Simulated Failures for Testing}
To test the fault tolerance mechanisms, we implemented simulated failures in the Crawler Nodes:

\begin{lstlisting}[language=Python, caption=Simulated Failure for Testing]
# Simulate occasional failures for testing fault tolerance
if random.random() < 0.05:  # 5% chance of simulated failure
    logging.warning(f"Simulating a brief node failure (testing fault tolerance)")
    time.sleep(12)  # Sleep longer than heartbeat timeout to trigger failure detection
\end{lstlisting}

\subsection{Monitoring System Implementation}
We have implemented a comprehensive monitoring system that tracks system metrics, node status, and performance data. The monitoring system consists of enhanced logging, centralized metrics collection, and a web-based monitoring dashboard.

\subsubsection{System Metrics Collection}
The Master Node collects and maintains system-wide metrics in a `system_metrics` dictionary:

\begin{lstlisting}[language=Python, caption=System Metrics Collection]
# Initialize system metrics in masterNode.py
system_metrics = {
    "start_time": datetime.now().isoformat(),
    "urls_crawled": 0,
    "urls_indexed": 0,
    "urls_queued": len(seed_urls),
    "urls_failed": 0,
    "error_count": 0,
    "crawler_status": {str(node): "active" for node in active_crawler_nodes},
    "crawler_performance": {str(node): {"urls_processed": 0, "avg_time_per_url": 0} for node in active_crawler_nodes},
    "indexer_status": "active",
    "task_assignments": [],
    "task_failures": []
}
\end{lstlisting}

These metrics are updated throughout the Master Node's operation:

\begin{lstlisting}[language=Python, caption=Updating System Metrics]
# Update metrics when processing extracted URLs
system_metrics["urls_crawled"] += 1
system_metrics["urls_queued"] += len(valid_new_urls)

# Update metrics for crawler performance
if str(source) in system_metrics["crawler_performance"]:
    perf = system_metrics["crawler_performance"][str(source)]
    perf["urls_processed"] += 1
    if task_start_times[source]:
        task_duration = (current_time - task_start_times[source]).total_seconds()
        # Update average time (running average)
        perf["avg_time_per_url"] = (perf["avg_time_per_url"] * (perf["urls_processed"] - 1) + task_duration) / perf["urls_processed"]

# Update metrics for indexing
if "Indexed URL" in message:
    system_metrics["urls_indexed"] += 1
\end{lstlisting}

\subsubsection{Monitoring Data Storage}
The system metrics are regularly written to a JSON file for persistence and access by the monitoring dashboard:

\begin{lstlisting}[language=Python, caption=Monitoring Data Storage]
def update_monitoring_data():
    """Update the monitoring data file with current system metrics"""
    # Ensure the monitoring directory exists
    monitoring_dir = os.path.join("data", "monitoring")
    os.makedirs(monitoring_dir, exist_ok=True)
    
    # Update last_updated timestamp
    system_metrics["last_updated"] = datetime.now().isoformat()
    
    # Write current metrics to the monitoring data file
    monitoring_file = os.path.join(monitoring_dir, "monitoring_data.json")
    with open(monitoring_file, "w") as f:
        json.dump(system_metrics, f, indent=2)
    
    logging.debug("Updated monitoring data file")
\end{lstlisting}

\subsubsection{Monitoring Dashboard}
We have implemented a Flask-based web application that provides real-time visualization of system metrics:

\begin{lstlisting}[language=Python, caption=Monitoring Dashboard Implementation]
app = Flask(__name__)
app.config['TEMPLATES_AUTO_RELOAD'] = True

# Initialize global data storage
system_data = {
    "last_update": datetime.now().isoformat(),
    "master_status": "Unknown",
    "crawlers": {},
    "indexer_status": "Unknown",
    "urls_crawled": 0,
    "urls_indexed": 0,
    "urls_failed": 0,
    "error_count": 0,
    "recent_tasks": []
}

def load_monitoring_data():
    """Load and parse the monitoring data JSON file"""
    global system_data
    try:
        monitoring_data_path = os.path.join("data", "monitoring", "monitoring_data.json")
        if os.path.exists(monitoring_data_path):
            with open(monitoring_data_path, "r") as f:
                data = json.load(f)
                
            # Update system_data with the loaded data
            system_data["last_update"] = datetime.now().isoformat()
            system_data["urls_crawled"] = data.get("urls_crawled", 0)
            system_data["urls_indexed"] = data.get("urls_indexed", 0)
            system_data["urls_failed"] = data.get("urls_failed", 0)
            system_data["error_count"] = data.get("error_count", 0)
            
            # Get crawler status
            system_data["crawlers"] = data.get("crawler_status", {})
            
            # Get crawler performance metrics
            system_data["crawler_performance"] = data.get("crawler_performance", {})
            
            # Get recent tasks (last 10)
            tasks = data.get("task_assignments", [])
            system_data["recent_tasks"] = sorted(
                tasks, 
                key=lambda x: x.get("time", ""), 
                reverse=True
            )[:10]
            
            # Calculate success rate
            total_urls = system_data["urls_crawled"] + system_data["urls_failed"]
            system_data["success_rate"] = (
                (system_data["urls_crawled"] / total_urls * 100) 
                if total_urls > 0 else 0
            )
            
            return True
    except Exception as e:
        print(f"Error loading monitoring data: {e}")
        return False
\end{lstlisting}

The dashboard provides API endpoints for accessing system metrics:

\begin{lstlisting}[language=Python, caption=Monitoring Dashboard API Endpoints]
@app.route('/api/data')
def get_data():
    """API endpoint to get the current system data"""
    return jsonify(system_data)

@app.route('/api/status')
def get_status():
    """API endpoint to get a summary of the system status"""
    active_crawlers = sum(1 for status in system_data["crawlers"].values() if status == "active")
    failed_crawlers = sum(1 for status in system_data["crawlers"].values() if status == "failed")
    
    status = {
        "master": system_data["master_status"],
        "active_crawlers": active_crawlers,
        "failed_crawlers": failed_crawlers,
        "total_crawlers": len(system_data["crawlers"]),
        "urls_crawled": system_data["urls_crawled"],
        "urls_indexed": system_data["urls_indexed"],
        "error_rate": f"{system_data.get('error_count', 0) / max(1, system_data['urls_crawled'] + system_data['urls_failed']) * 100:.1f}%"
    }
    
    return jsonify(status)
\end{lstlisting}

The dashboard updates its data every 2 seconds to provide near real-time monitoring:

\begin{lstlisting}[language=Python, caption=Background Data Update]
def background_data_update():
    """Background thread to update data periodically"""
    while True:
        load_monitoring_data()
        read_log_files()
        time.sleep(2)  # Update every 2 seconds

# Start the background update thread
update_thread = threading.Thread(target=background_data_update, daemon=True)
update_thread.start()
\end{lstlisting} 

\subsection{Cloud Storage Implementation}
To ensure data durability and enable potential re-indexing, we have implemented cloud storage for both raw HTML and processed text data using AWS S3. The cloud storage system includes fallback to local storage when cloud services are unavailable.

\subsubsection{AWS S3 Integration}
The `CloudStorage` class provides a robust interface for storing and retrieving data from AWS S3:

\begin{lstlisting}[language=Python, caption=CloudStorage Initialization]
def __init__(self, bucket_name=None, region_name=None):
    """
    Initialize the CloudStorage class.
    
    Args:
        bucket_name (str): The S3 bucket name to use. If None, will use environment variable or default.
        region_name (str): The AWS region to use.
    """
    self.logger = logging.getLogger('CloudStorage')
    self.logger.setLevel(logging.INFO)
    
    # Try to load AWS configuration from file
    try:
        if os.path.exists('aws_config.json'):
            with open('aws_config.json', 'r') as f:
                aws_config = json.load(f)
                self.logger.info("Loaded AWS configuration from aws_config.json")
                config_bucket = aws_config.get('aws', {}).get('bucket_name')
                config_region = aws_config.get('aws', {}).get('region')
                
                # Use provided values or config values
                bucket_name = bucket_name or config_bucket
                region_name = region_name or config_region
    except Exception as e:
        self.logger.warning(f"Failed to load AWS configuration file: {e}")
    
    # Get bucket name from environment variable if not provided
    self.bucket_name = bucket_name or os.environ.get('AWS_S3_BUCKET', 'web-crawler-data-storage')
    self.region_name = region_name or os.environ.get('AWS_DEFAULT_REGION', 'us-east-1')
    
    # Initialize AWS clients
    try:
        self.s3_client = boto3.client('s3', region_name=self.region_name)
        self.s3_resource = boto3.resource('s3', region_name=self.region_name)
        self._ensure_bucket_exists()
        self.logger.info(f"Successfully connected to AWS S3 in region {self.region_name}")
    except Exception as e:
        self.logger.error(f"Failed to initialize S3 client: {e}")
        # Initialize to None to allow graceful fallback to local storage
        self.s3_client = None
        self.s3_resource = None
\end{lstlisting}

\subsubsection{Storing Raw HTML and Processed Text}
The `CloudStorage` class provides methods for storing both raw HTML and processed text data:

\begin{lstlisting}[language=Python, caption=CloudStorage Data Storage]
def store_raw_html(self, url, html_content):
    """
    Store raw HTML content in S3.
    
    Args:
        url (str): The URL of the crawled page
        html_content (str): The raw HTML content
        
    Returns:
        dict: Storage info including success status and storage location
    """
    if not self.s3_client:
        self.logger.warning("S3 client not initialized. Falling back to local storage.")
        return self._local_store(url, html_content, 'raw_html')
    
    key = self._generate_key(url, 'raw_html')
    try:
        self.s3_client.put_object(
            Bucket=self.bucket_name,
            Key=key,
            Body=html_content,
            ContentType='text/html',
            Metadata={
                'url': url,
                'timestamp': str(int(time.time())),
                'storage_type': 'raw_html'
            }
        )
        self.logger.info(f"Successfully stored raw HTML for {url} at s3://{self.bucket_name}/{key}")
        return {
            'success': True,
            'storage_type': 's3',
            'bucket': self.bucket_name,
            'key': key,
            'url': url,
            'content_type': 'raw_html',
            'timestamp': datetime.now().isoformat()
        }
    except Exception as e:
        self.logger.error(f"Failed to store raw HTML for {url}: {e}")
        # Fall back to local storage
        return self._local_store(url, html_content, 'raw_html')
\end{lstlisting}

\subsubsection{Local Storage Fallback}
The system includes a local storage fallback mechanism to ensure data is preserved even when cloud storage is unavailable:

\begin{lstlisting}[language=Python, caption=Local Storage Fallback]
def _local_store(self, url, content, content_type, metadata=None):
    """
    Store content locally when cloud storage is unavailable.
    
    Args:
        url (str): The URL of the crawled page
        content (str): The content to store
        content_type (str): Type of content (raw_html, processed_text, metadata)
        metadata (dict): Optional metadata about the content
        
    Returns:
        dict: Storage info including success status and storage location
    """
    try:
        # Create a hash of the URL
        url_hash = hashlib.md5(url.encode()).hexdigest()
        
        # Create directory structure
        local_dir = os.path.join("crawled_data", content_type)
        os.makedirs(local_dir, exist_ok=True)
        
        # Determine file extension
        extension = 'html' if content_type == 'raw_html' else 'json' if content_type == 'metadata' else 'txt'
        
        # Create filename
        filename = f"{url_hash}.{extension}"
        filepath = os.path.join(local_dir, filename)
        
        # Write content to file
        with open(filepath, 'w', encoding='utf-8') as f:
            if content_type == 'metadata' and metadata:
                json.dump(metadata, f, indent=2)
            else:
                f.write(content)
        
        self.logger.info(f"Successfully stored {content_type} for {url} locally at {filepath}")
        
        # If metadata is provided separately, store it as well
        if metadata and content_type != 'metadata':
            metadata_path = os.path.join("crawled_data", "metadata", f"{url_hash}.json")
            os.makedirs(os.path.dirname(metadata_path), exist_ok=True)
            with open(metadata_path, 'w', encoding='utf-8') as f:
                json.dump(metadata, f, indent=2)
        
        return {
            'success': True,
            'storage_type': 'local',
            'filepath': filepath,
            'url': url,
            'content_type': content_type,
            'timestamp': datetime.now().isoformat()
        }
    except Exception as e:
        self.logger.error(f"Failed to store {content_type} locally for {url}: {e}")
        return {
            'success': False,
            'storage_type': 'none',
            'error': str(e)
        }
\end{lstlisting}

\subsubsection{Integration with the Indexer}
The `CloudStorage` class is integrated with the Indexer Node to store both raw HTML and processed text during indexing:

\begin{lstlisting}[language=Python, caption=CloudStorage Integration with Indexer]
# Store the raw HTML in cloud storage if available
if self.cloud_storage:
    raw_html_result = self.cloud_storage.store_raw_html(url, content)
    logging.info(f"Stored raw HTML for {url} with result: {raw_html_result['storage_type']}")
    
    # Store processed text and metadata
    processed_result = self.cloud_storage.store_processed_text(url, extracted_text, metadata)
    logging.info(f"Stored processed text for {url} with result: {processed_result['storage_type']}")
\end{lstlisting}

\subsection{Search Interface Implementation}
We have implemented a web-based search interface that allows users to submit search queries with advanced features such as boolean operators and phrase search. The search interface leverages the enhanced indexing capabilities implemented in Phase 3.

\subsubsection{Search API Implementation}
The search API is implemented as a Flask web application that provides an interface to the Whoosh index:

\begin{lstlisting}[language=Python, caption=Search API Implementation]
from flask import Flask, render_template, request, jsonify
from whoosh.index import open_dir
from whoosh.qparser import MultifieldParser
from whoosh import scoring
import os

app = Flask(__name__)

# Initialize the search index
INDEX_DIR = "search_index"

def get_searcher():
    if not os.path.exists(INDEX_DIR):
        return None
    return open_dir(INDEX_DIR).searcher(weighting=scoring.BM25F)

@app.route('/search')
def search():
    query = request.args.get('q', '')
    page = int(request.args.get('page', 1))
    
    if not query:
        return jsonify({'results': [], 'total': 0, 'page': page})
    
    try:
        with get_searcher() as searcher:
            if searcher is None:
                return jsonify({'error': 'Search index not found. Please run the crawler first.'})
                
            schema = searcher.schema
            fields = ["title", "content", "keywords"]
            parser = MultifieldParser(fields, schema)
            q = parser.parse(query)
            
            # Search with pagination (10 results per page)
            results = searcher.search_page(q, page, pagelen=10)
            
            response = {
                'results': [
                    {
                        'title': hit['title'] if hit['title'] else hit['url'],
                        'url': hit['url'],
                        'summary': hit['summary'],
                        'keywords': hit['keywords'],
                        'score': hit.score,
                        'last_updated': hit['last_updated']
                    } for hit in results
                ],
                'total': len(results),
                'page': page,
                'total_pages': (len(results) + 9) // 10
            }
            
            return jsonify(response)
            
    except Exception as e:
        return jsonify({'error': str(e)})
\end{lstlisting}

\subsubsection{Search Interface UI}
The search interface provides a user-friendly web UI for submitting search queries and viewing results:

\begin{lstlisting}[language=HTML, caption=Search Interface HTML Template]
<!DOCTYPE html>
<html lang="en">
<head>
    <meta charset="UTF-8">
    <meta name="viewport" content="width=device-width, initial-scale=1.0">
    <title>Web Crawler Search Interface</title>
    <link href="https://cdn.jsdelivr.net/npm/bootstrap@5.1.3/dist/css/bootstrap.min.css" rel="stylesheet">
    <style>
        .search-result {
            margin-bottom: 2rem;
            padding: 1rem;
            border-bottom: 1px solid #eee;
        }
        .search-result:hover {
            background-color: #f8f9fa;
        }
        .score-badge {
            font-size: 0.8rem;
            color: #6c757d;
        }
        .keywords {
            font-size: 0.9rem;
            color: #28a745;
        }
        .last-updated {
            font-size: 0.8rem;
            color: #6c757d;
        }
        #searchHelp {
            font-size: 0.9rem;
            color: #6c757d;
            margin-top: 0.5rem;
        }
    </style>
</head>
<body>
    <div class="container mt-5">
        <h1 class="text-center mb-4">Web Crawler Search</h1>
        
        <div class="row justify-content-center">
            <div class="col-md-8">
                <div class="input-group mb-3">
                    <input type="text" id="searchInput" class="form-control" placeholder="Enter your search query...">
                    <button class="btn btn-primary" onclick="search(1)">Search</button>
                </div>
                <div id="searchHelp">
                    <strong>Search Tips:</strong><br>
                    • Use AND, OR, NOT for boolean search (e.g., "python AND programming")<br>
                    • Use quotes for exact phrases (e.g., "machine learning")<br>
                    • Use field-specific search (e.g., title:github, content:"open source")<br>
                </div>
            </div>
        </div>

        <div id="results" class="mt-4">
            <!-- Search results will be displayed here -->
        </div>

        <div id="pagination" class="d-flex justify-content-center mt-4">
            <!-- Pagination will be displayed here -->
        </div>
    </div>

    <script>
        function search(page) {
            const query = document.getElementById('searchInput').value;
            if (!query) return;

            fetch(`/search?q=${encodeURIComponent(query)}&page=${page}`)
                .then(response => response.json())
                .then(data => {
                    // Display search results
                })
                .catch(error => {
                    document.getElementById('results').innerHTML = 
                      `<div class="alert alert-danger">Error: ${error}</div>`;
                });
        }
    </script>
</body>
</html>
\end{lstlisting}

\section{Testing and Evaluation}

\subsection{Fault Tolerance Testing}
We conducted extensive testing of the fault tolerance mechanisms to ensure the system can handle Crawler Node failures gracefully. The testing included both simulated failures and real-world failure scenarios.

\subsubsection{Simulated Failure Testing}
To test the fault tolerance mechanisms, we implemented simulated failures in the Crawler Nodes with a 5\% probability of failure:

\begin{lstlisting}[language=Python, caption=Simulated Failure Code]
# Simulate occasional failures for testing fault tolerance
if random.random() < 0.05:  # 5% chance of simulated failure
    logging.warning(f"Simulating a brief node failure (testing fault tolerance)")
    time.sleep(12)  # Sleep longer than heartbeat timeout to trigger failure detection
\end{lstlisting}

The simulated failures allowed us to test the heartbeat monitoring, task timeout detection, and task re-queueing mechanisms in a controlled environment.

\subsubsection{Test Results}
The system successfully detected and handled both simulated and real Crawler Node failures:

\begin{enumerate}
    \item \textbf{Heartbeat Monitoring}: The system correctly detected when a Crawler Node failed to send heartbeats and marked it as failed.
    \item \textbf{Task Timeout Detection}: The system correctly identified tasks that exceeded the timeout threshold and re-queued them for processing.
    \item \textbf{Task Re-queueing}: Failed tasks were successfully re-queued and assigned to other active Crawler Nodes.
    \item \textbf{Recovery Detection}: When failed nodes became active again, the system correctly detected this and updated their status.
\end{enumerate}

\subsection{Indexing and Search Testing}
We conducted thorough testing of the enhanced indexing and search functionality to ensure it meets the Phase 3 requirements.

\subsubsection{Indexing Tests}
The indexing tests focused on evaluating the text processing, keyword extraction, and summarization capabilities:

\begin{enumerate}
    \item \textbf{Text Processing}: The system correctly processed HTML content, extracted text, removed stopwords, and applied lemmatization.
    \item \textbf{Keyword Extraction}: The system successfully extracted relevant keywords from the processed text.
    \item \textbf{Summarization}: The system generated concise summaries of the indexed content.
    \item \textbf{Index Performance}: The Whoosh-based index demonstrated efficient storage and retrieval capabilities.
\end{enumerate}

\subsubsection{Search Tests}
The search tests evaluated the advanced search capabilities:

\begin{enumerate}
    \item \textbf{Boolean Operators}: The search engine correctly handled boolean operators (AND, OR, NOT) in queries.
    \item \textbf{Phrase Search}: The search engine correctly processed phrase queries enclosed in quotes.
    \item \textbf{Field-Specific Search}: The search engine correctly handled field-specific queries using the field:term syntax.
    \item \textbf{Relevance Ranking}: The BM25F scoring algorithm provided relevant result ranking.
\end{enumerate}

\subsection{Cloud Storage Testing}
We tested the cloud storage functionality to ensure data durability and availability:

\begin{enumerate}
    \item \textbf{S3 Storage}: The system successfully stored both raw HTML and processed text in AWS S3.
    \item \textbf{Local Fallback}: When AWS S3 was unavailable, the system correctly fell back to local storage.
    \item \textbf{Data Retrieval}: The system could retrieve stored data from both S3 and local storage.
    \item \textbf{Metadata Storage}: The system correctly stored and retrieved metadata about the indexed content.
\end{enumerate}

\subsection{Scalability Testing}
We conducted basic scalability testing by running the system with varying numbers of Crawler Nodes:

\begin{enumerate}
    \item \textbf{Single Crawler}: The system operated correctly with a single Crawler Node.
    \item \textbf{Multiple Crawlers}: The system demonstrated improved throughput with multiple Crawler Nodes.
    \item \textbf{Task Distribution}: The Master Node correctly distributed tasks among available Crawler Nodes.
    \item \textbf{Load Balancing}: The system balanced the load among active Crawler Nodes.
\end{enumerate}

\section{Implementation Evidence}
This section provides evidence from our system logs that demonstrates the successful implementation of the Phase 3 requirements. These log captures show the system in operation, highlighting the key features implemented.

\subsection{Enhanced Indexing Evidence}
The following log extract from the Indexer Node shows the Whoosh indexing in action, including advanced text processing, keyword extraction, and cloud storage integration:

\begin{lstlisting}[language=bash, caption=Indexer Log - Advanced Indexing and Keyword Extraction]
2025-04-18 00:07:32,748 - Indexer - INFO - Successfully indexed https://www.python.org
2025-04-18 00:07:32,748 - Indexer - INFO - Extracted 10 keywords
2025-04-18 00:07:33,264 - Indexer - INFO - Successfully indexed https://www.reddit.com
2025-04-18 00:07:33,264 - Indexer - INFO - Extracted 10 keywords
2025-04-18 00:07:33,613 - Indexer - INFO - Successfully indexed https://www.wikipedia.org
2025-04-18 00:07:33,614 - Indexer - INFO - Extracted 10 keywords
2025-04-18 00:07:34,156 - Indexer - INFO - Successfully indexed https://www.github.com
2025-04-18 00:07:34,156 - Indexer - INFO - Extracted 10 keywords
2025-04-18 00:07:47,173 - Indexer - INFO - Indexer received shutdown signal. Generating final statistics...
2025-04-18 00:07:47,205 - Indexer - INFO - Final Index Statistics:
2025-04-18 00:07:47,206 - Indexer - INFO - Total documents indexed: 23
2025-04-18 00:07:47,206 - Indexer - INFO - Total unique URLs processed: 20
2025-04-18 00:07:47,206 - Indexer - INFO - 
Example Search Results:
2025-04-18 00:07:47,248 - Indexer - INFO - 
Query: python AND programming
2025-04-18 00:07:47,248 - Indexer - INFO - Found 5 results
2025-04-18 00:07:47,253 - Indexer - INFO - 
Query: title:github
2025-04-18 00:07:47,254 - Indexer - INFO - Found 0 results
2025-04-18 00:07:47,274 - Indexer - INFO - 
Query: "open source"
2025-04-18 00:07:47,274 - Indexer - INFO - Found 9 results
2025-04-18 00:07:47,294 - Indexer - INFO - 
Query: content:machine learning
2025-04-18 00:07:47,294 - Indexer - INFO - Found 3 results
\end{lstlisting}

The logs show the following advanced features:
\begin{itemize}
    \item Automatic keyword extraction for each indexed document
    \item Support for boolean search operators (AND) in search queries
    \item Field-specific search using field prefix notation (title:github)
    \item Phrase search using quotes ("open source")
    \item Content-specific field searches (content:machine learning)
\end{itemize}

\subsection{Cloud Storage Implementation Evidence}
The system implements robust cloud storage with AWS S3 integration, as demonstrated by these logs:

\begin{lstlisting}[language=bash, caption=Cloud Storage Initialization and Configuration]
2025-04-18 23:08:46,083 - CloudStorage - INFO - Loaded AWS configuration from aws_config.json
2025-04-18 23:08:47,276 - CloudStorage - INFO - Bucket web-crawler-data-storage already exists
2025-04-18 23:08:47,277 - CloudStorage - INFO - Successfully connected to AWS S3 in region us-east-1
\end{lstlisting}

The log entries above show the successful initialization of the CloudStorage component, including:
\begin{itemize}
    \item Loading AWS configuration from an external configuration file
    \item Verifying the existence of the designated S3 bucket
    \item Establishing a connection to AWS S3 in the specified region
\end{itemize}

The following logs demonstrate successful storage of various data types in AWS S3:

\begin{lstlisting}[language=bash, caption=Successful Storage of Multiple Data Types in AWS S3]
2025-04-18 00:44:19,468 - CloudStorage - INFO - Successfully stored raw HTML for https://www.python.org at s3://web-crawler-data-storage/raw_html/2025/04/18/c137682bbb0946674f25d1d4bd9e07a0.html
2025-04-18 00:44:19,773 - CloudStorage - INFO - Successfully stored processed text for https://www.python.org at s3://web-crawler-data-storage/processed_text/2025/04/18/c137682bbb0946674f25d1d4bd9e07a0.txt
2025-04-18 00:44:20,076 - CloudStorage - INFO - Successfully stored metadata for https://www.python.org at s3://web-crawler-data-storage/metadata/2025/04/18/c137682bbb0946674f25d1d4bd9e07a0.json

2025-04-18 00:44:21,607 - CloudStorage - INFO - Successfully stored raw HTML for https://www.github.com at s3://web-crawler-data-storage/raw_html/2025/04/18/1895ab41ad3a81f6fd0a194bbc337a22.html
2025-04-18 00:44:21,925 - CloudStorage - INFO - Successfully stored processed text for https://www.github.com at s3://web-crawler-data-storage/processed_text/2025/04/18/1895ab41ad3a81f6fd0a194bbc337a22.txt
2025-04-18 00:44:22,678 - CloudStorage - INFO - Successfully stored metadata for https://www.github.com at s3://web-crawler-data-storage/metadata/2025/04/18/1895ab41ad3a81f6fd0a194bbc337a22.json
\end{lstlisting}

These storage operations demonstrate several key features of our cloud storage implementation:
\begin{itemize}
    \item Hierarchical organization of data in S3 with separate paths for raw HTML, processed text, and metadata
    \item Date-based partitioning (year/month/day) for efficient data organization
    \item Content-specific file extensions (.html, .txt, .json) for appropriate data type identification
    \item Consistent hashing of URLs to generate unique filenames
    \item Complete storage workflow for each crawled page, ensuring all three data types are preserved
\end{itemize}

The logs also show that our Master node receives confirmation when content is successfully indexed with cloud storage:

\begin{lstlisting}[language=bash, caption=Master Node Confirmation of Cloud Storage Integration]
2025-04-18 00:44:15,785 - Master - INFO - Status from node 4: Indexed https://www.wikipedia.org with 10 keywords with cloud storage
2025-04-18 00:44:18,913 - Master - INFO - Status from node 4: Indexed https://www.reddit.com with 10 keywords with cloud storage
2025-04-18 00:44:20,221 - Master - INFO - Status from node 4: Indexed https://www.python.org with 10 keywords with cloud storage
2025-04-18 00:44:23,556 - Master - INFO - Status from node 4: Indexed https://www.github.com with 10 keywords with cloud storage
\end{lstlisting}

\subsection{Fault Tolerance Evidence}
The following logs provide comprehensive evidence of our fault tolerance mechanisms in action, including heartbeat monitoring, failure detection, and automatic task reassignment:

\begin{lstlisting}[language=bash, caption=Heartbeat Monitoring from Multiple Crawler Nodes]
2025-04-18 00:44:07,629 - Master - INFO - Status from node 2: Heartbeat from crawler 2: Active
2025-04-18 00:44:07,832 - Master - INFO - Status from node 3: Heartbeat from crawler 3: Active
2025-04-18 00:44:08,033 - Master - INFO - Status from node 1: Heartbeat from crawler 1: Active
\end{lstlisting}

A critical failure detection and recovery scenario is captured in the following log sequence:

\begin{lstlisting}[language=bash, caption=Failure Detection and Recovery]
2025-04-18 00:44:51,071 - Master - WARNING - Heartbeat timeout for crawler 1. Marking as failed.
2025-04-18 00:44:51,072 - Master - WARNING - No available crawler nodes found. Waiting for nodes to become available.
2025-04-18 00:44:51,172 - Master - WARNING - No available crawler nodes found. Waiting for nodes to become available.
2025-04-18 00:44:51,273 - Master - WARNING - No available crawler nodes found. Waiting for nodes to become available.
2025-04-18 00:44:51,374 - Master - INFO - Status from node 3: Heartbeat from crawler 3: Active
2025-04-18 00:44:51,375 - Master - WARNING - No available crawler nodes found. Waiting for nodes to become available.
...
2025-04-18 00:44:52,182 - Master - INFO - Crawler 1 is back online.
2025-04-18 00:44:52,182 - Master - INFO - Status from node 1: Heartbeat from crawler 1: Active
2025-04-18 00:44:52,183 - Master - INFO - Assigned task 16 to crawler 1 for URL: https://github.com/security/advanced-security
\end{lstlisting}

This sequence demonstrates several key fault tolerance features:
\begin{itemize}
    \item Detection of a missed heartbeat from crawler 1
    \item Immediate marking of the node as failed
    \item Continuous monitoring for available nodes
    \item Quick detection of crawler 1's recovery (approximately 1.1 seconds after failure detection)
    \item Immediate task assignment to the recovered node
\end{itemize}

The following logs from crawler nodes show the simulated failures that were used to test the fault tolerance mechanisms:

\begin{lstlisting}[language=bash, caption=Simulated Failures from Multiple Crawler Nodes]
# Crawler 1 Simulated Failure:
2025-04-18 00:44:52,132 - Crawler-41624 - INFO - Sent content to indexer node for https://pypi.org/
2025-04-18 00:44:52,132 - Crawler-41624 - WARNING - Simulating a brief node failure (testing fault tolerance)
2025-04-18 00:45:04,133 - Crawler-41624 - INFO - Task for https://pypi.org/ completed in 24.65 seconds

# Crawler 2 Simulated Failure:
2025-04-18 00:42:12,736 - Crawler - INFO - Crawler 2 completed task for URL: https://www.github.com (found 68 URLs in 3.82 seconds)
2025-04-18 00:42:14,139 - Crawler - INFO - Crawler 2 received URL: https://meta.wikimedia.org/wiki/Special:MyLanguage/List_of_Wikipedias
2025-04-18 00:43:22,834 - Crawler - WARNING - Simulating a brief node failure (testing fault tolerance)
2025-04-18 00:43:34,986 - Crawler - INFO - Crawler 2 completed task for URL: https://www.reddit.com/login/ (found 22 URLs in 5.44 seconds)
\end{lstlisting}

The Master node's response to these failures shows the complete fault tolerance workflow:
1. Regular heartbeat monitoring:
\begin{lstlisting}[language=bash]
2025-04-18 00:44:07,629 - Master - INFO - Status from node 2: Heartbeat from crawler 2: Active
2025-04-18 00:44:07,832 - Master - INFO - Status from node 3: Heartbeat from crawler 3: Active
2025-04-18 00:44:08,033 - Master - INFO - Status from node 1: Heartbeat from crawler 1: Active
\end{lstlisting}

2. Failure detection:
\begin{lstlisting}[language=bash]
2025-04-18 00:44:51,071 - Master - WARNING - Heartbeat timeout for crawler 1. Marking as failed.
\end{lstlisting}

3. Recovery detection:
\begin{lstlisting}[language=bash]
2025-04-18 00:44:52,182 - Master - INFO - Crawler 1 is back online.
2025-04-18 00:44:52,182 - Master - INFO - Status from node 1: Heartbeat from crawler 1: Active
\end{lstlisting}

4. Task reassignment:
\begin{lstlisting}[language=bash]
2025-04-18 00:44:52,183 - Master - INFO - Assigned task 16 to crawler 1 for URL: https://github.com/security/advanced-security
\end{lstlisting}

The system's fault tolerance capabilities were rigorously tested using simulated failures across multiple crawler nodes, demonstrating robust failure detection, recovery monitoring, and efficient task reassignment to ensure continuous system operation.

\subsection{Performance and Task Completion Evidence}
The following logs demonstrate the system's task assignment, processing, and monitoring capabilities:

\begin{lstlisting}[language=bash, caption=Crawler Performance Metrics]
2025-04-18 00:42:55,232 - Crawler - INFO - Crawler 2 fetching URL: https://psfmember.org/civicrm/contribute/transact?reset=1&id=2
2025-04-18 00:42:58,138 - Crawler - INFO - Crawler 2 completed task for URL: https://psfmember.org/civicrm/contribute/transact?reset=1&id=2 (found 10 URLs in 2.91 seconds)
2025-04-18 00:42:58,140 - Crawler - INFO - Sending content to indexer node
2025-04-18 00:42:58,140 - Crawler - INFO - Sending extracted URLs to master
2025-04-18 00:42:59,642 - Crawler - INFO - Crawler 2 received URL: https://github.com/features/copilot
2025-04-18 00:43:04,245 - Crawler - INFO - Crawler 2 completed task for URL: https://github.com/features/copilot (found 122 URLs in 4.60 seconds)
\end{lstlisting}

\begin{lstlisting}[language=bash, caption=Master Log - Task Assignment and Monitoring]
2025-04-18 00:44:05,216 - Master - INFO - Master node started with rank 0 of 5
2025-04-18 00:44:05,217 - Master - INFO - Active Crawler Nodes: [1, 2, 3]
2025-04-18 00:44:05,217 - Master - INFO - Indexer Node: 4
2025-04-18 00:44:05,218 - Master - INFO - Assigned task 0 to crawler 1 for URL: https://www.python.org
2025-04-18 00:44:05,220 - Master - INFO - Assigned task 1 to crawler 2 for URL: https://www.github.com
2025-04-18 00:44:05,222 - Master - INFO - Assigned task 2 to crawler 3 for URL: https://www.wikipedia.org
2025-04-18 00:44:07,226 - Master - INFO - Assigned task 2 to crawler 3 for URL: https://www.wikipedia.org
2025-04-18 00:44:07,629 - Master - INFO - Status from node 2: Heartbeat from crawler 2: Active
2025-04-18 00:44:07,832 - Master - INFO - Status from node 3: Heartbeat from crawler 3: Active
2025-04-18 00:44:08,033 - Master - INFO - Status from node 1: Heartbeat from crawler 1: Active
2025-04-18 00:44:08,235 - Master - INFO - Received URLs from crawler 3. Queue size now: 7, Processed URLs: 1
2025-04-18 00:44:08,237 - Master - INFO - Assigned task 3 to crawler 3 for URL: https://www.reddit.com
\end{lstlisting}

\subsection{Comprehensive System Operation}
The indexer logs demonstrate the complete system workflow, including advanced text processing, keyword extraction, and successful data storage:

\begin{lstlisting}[language=bash, caption=Indexer Log - Complete Processing Workflow]
2025-04-18 00:07:32,748 - Indexer - INFO - Successfully indexed https://www.python.org
2025-04-18 00:07:32,748 - Indexer - INFO - Extracted 10 keywords
2025-04-18 00:07:33,127 - CloudStorage - WARNING - S3 client not initialized. Falling back to local storage.
2025-04-18 00:07:33,135 - CloudStorage - INFO - Successfully stored raw_html for https://www.reddit.com locally at crawled_data\raw_html\30dd8fd5626b7e36783516a6e238baf2.html
2025-04-18 00:07:33,136 - CloudStorage - WARNING - S3 client not initialized. Falling back to local storage.
2025-04-18 00:07:33,139 - CloudStorage - INFO - Successfully stored processed_text for https://www.reddit.com locally at crawled_data\processed_text\30dd8fd5626b7e36783516a6e238baf2.txt
2025-04-18 00:07:33,264 - Indexer - INFO - Successfully indexed https://www.reddit.com
2025-04-18 00:07:33,264 - Indexer - INFO - Extracted 10 keywords

2025-04-18 00:38:37,541 - CloudStorage - INFO - Successfully stored raw HTML for https://www.stackoverflow.com at s3://web-crawler-data-storage/raw_html/2025/04/18/2415620324baad2a5308e2c497763b12.html
2025-04-18 00:38:37,902 - CloudStorage - INFO - Successfully stored processed text for https://www.stackoverflow.com at s3://web-crawler-data-storage/processed_text/2025/04/18/2415620324baad2a5308e2c497763b12.txt
2025-04-18 00:38:38,260 - CloudStorage - INFO - Successfully stored metadata for https://www.stackoverflow.com at s3://web-crawler-data-storage/metadata/2025/04/18/2415620324baad2a5308e2c497763b12.json
\end{lstlisting}

\begin{lstlisting}[language=bash, caption=Master Node - Final Statistics]
2025-04-18 00:07:47,173 - Indexer - INFO - Indexer received shutdown signal. Generating final statistics...
2025-04-18 00:07:47,205 - Indexer - INFO - Final Index Statistics:
2025-04-18 00:07:47,206 - Indexer - INFO - Total documents indexed: 23
2025-04-18 00:07:47,206 - Indexer - INFO - Total unique URLs processed: 20
2025-04-18 00:07:47,206 - Indexer - INFO - 
Example Search Results:
2025-04-18 00:07:47,248 - Indexer - INFO - 
Query: python AND programming
2025-04-18 00:07:47,248 - Indexer - INFO - Found 5 results
2025-04-18 00:07:47,253 - Indexer - INFO - 
Query: title:github
2025-04-18 00:07:47,254 - Indexer - INFO - Found 0 results
2025-04-18 00:07:47,274 - Indexer - INFO - 
Query: "open source"
2025-04-18 00:07:47,274 - Indexer - INFO - Found 9 results
2025-04-18 00:07:47,294 - Indexer - INFO - 
Query: content:machine learning
2025-04-18 00:07:47,294 - Indexer - INFO - Found 3 results
\end{lstlisting}

These log captures provide comprehensive evidence that all the required Phase 3 features have been successfully implemented in our system, including:

\begin{enumerate}
    \item \textbf{Enhanced Indexing}: Advanced text processing with keyword extraction and support for complex search queries
    \item \textbf{Cloud Storage}: Integration with AWS S3 for durable storage of HTML, processed text, and metadata
    \item \textbf{Fault Tolerance}: Robust heartbeat monitoring, failure detection, and automatic recovery
    \item \textbf{Local Fallback}: Graceful degradation to local storage when cloud services are unavailable
    \item \textbf{Task Coordination}: Efficient distribution and monitoring of crawling tasks
    \item \textbf{Performance Tracking}: Detailed metrics for system performance analysis
    \item \textbf{Monitoring Dashboard}: Real-time visualization of system status and metrics
\end{enumerate}

The system has demonstrated its ability to handle node failures gracefully, store data reliably, process content with advanced techniques, and provide powerful search capabilities. These features make our distributed web crawling and indexing system robust, reliable, and scalable.

\section{Conclusion and Future Work}
\subsection{Summary of Phase 3 Achievements}
In Phase 3, we have successfully enhanced our distributed web crawling and indexing system with the following improvements:

\begin{enumerate}
    \item \textbf{Enhanced Indexing}: We implemented a robust indexing system using Whoosh with advanced text processing, keyword extraction, and summarization capabilities.
    \item \textbf{Fault Tolerance}: We implemented comprehensive fault tolerance mechanisms including heartbeat monitoring, task timeout detection, and task re-queueing.
    \item \textbf{Monitoring System}: We developed a monitoring system with enhanced logging, centralized metrics collection, and a web-based dashboard.
    \item \textbf{Data Persistence}: We implemented cloud storage using AWS S3 for both raw HTML and processed text with local fallback capabilities.
    \item \textbf{Search Interface}: We created a web-based search interface with advanced search capabilities including boolean operators and phrase search.
\end{enumerate}

\subsection{Lessons Learned}
During the implementation of Phase 3, we learned several important lessons:

\begin{enumerate}
    \item \textbf{Fault Tolerance Design}: Designing and implementing fault tolerance mechanisms requires careful consideration of failure scenarios and recovery strategies.
    \item \textbf{Monitoring Importance}: A comprehensive monitoring system is essential for understanding system behavior, identifying issues, and ensuring optimal performance.
    \item \textbf{Cloud Integration}: Integrating with cloud services requires robust error handling and fallback mechanisms to handle connectivity issues.
    \item \textbf{Advanced Indexing}: Implementing advanced indexing capabilities requires a good understanding of NLP techniques and search algorithms.
    \item \textbf{User Interface Design}: Designing an effective search interface requires balancing simplicity with advanced features.
\end{enumerate}

\section{User Interface Documentation}

\subsection{Monitoring Dashboard Interface}
The monitoring dashboard provides a real-time graphical interface to visualize the system's operational status and performance metrics. The dashboard is implemented using Flask and Bootstrap, offering a responsive and modern user interface.

\subsubsection{Dashboard Layout and Components}
The monitoring dashboard features a clean, card-based layout with the following main components:

\begin{enumerate}
    \item \textbf{System Status Overview Cards}: 
    \begin{itemize}
        \item \textbf{Master Node Status}: Displays the current state of the Master Node (Running/Stopped)
        \item \textbf{Crawler Nodes}: Shows the count of active and failed crawler nodes with color-coded indicators
        \item \textbf{URLs Processed}: Displays the total number of URLs successfully crawled
        \item \textbf{Error Rate}: Shows the percentage of failed crawling attempts with error count
    \end{itemize}
    
    \item \textbf{Data Visualization Charts}:
    \begin{itemize}
        \item \textbf{Crawling Progress}: A doughnut chart that visualizes the proportion of crawled, indexed, and failed URLs
        \item \textbf{System Metrics}: A bar chart showing key metrics including active crawlers, URLs crawled, URLs indexed, and errors
    \end{itemize}
    
    \item \textbf{Crawler Node Status Table}: A detailed table showing:
    \begin{itemize}
        \item Crawler ID
        \item Current status (active/failed)
        \item URLs assigned and completed
        \item Failure rate for each crawler
    \end{itemize}
    
    \item \textbf{Recent Tasks Table}: Lists the most recent crawling tasks with:
    \begin{itemize}
        \item Timestamp
        \item URL being processed
        \item Assigned crawler node
        \item Current status (assigned/completed)
    \end{itemize}
    
    \item \textbf{Log Viewers}:
    \begin{itemize}
        \item \textbf{Master Node Logs}: Real-time log display from the Master Node
        \item \textbf{Crawler Node Logs}: Real-time log display from the Crawler Nodes
    \end{itemize}
\end{enumerate}

\begin{figure}[h]
\centering
\includegraphics[width=\textwidth]{dashboard_overview.png}
\caption{Web Crawler Monitoring Dashboard Overview}
\label{fig:dashboard_overview}
\end{figure}

\subsubsection{Key Metrics Visualization}
The dashboard visualizes several key performance indicators:

\begin{itemize}
    \item \textbf{Master Node Status}: Indicated with a green "Running" badge when operational
    \item \textbf{Crawler Nodes}: Displays "3 Active" and "0 Failed" with corresponding green and red badges
    \item \textbf{URLs Processed}: Shows the count of successfully crawled URLs (increasing from 6 to 16 to 20 as shown in the screenshots)
    \item \textbf{Error Rate}: Maintains "0.0\%" with "0 Errors" indicating reliable system operation
\end{itemize}

\begin{figure}[h]
\centering
\includegraphics[width=\textwidth]{dashboard_metrics.png}
\caption{Dashboard Key Metrics and Performance Indicators}
\label{fig:dashboard_metrics}
\end{figure}

\subsubsection{Dynamic Data Updates}
The dashboard updates automatically every 2 seconds to provide real-time visibility into the system's state:

\begin{itemize}
    \item System status changes are reflected immediately with appropriate color coding
    \item Task assignments and completions are logged in the Recent Tasks table
    \item The crawling progress chart updates to show the current distribution of crawled, indexed, and failed URLs
    \item The Master and Crawler logs scroll to show the most recent activities
\end{itemize}

\begin{figure}[h]
\centering
\includegraphics[width=\textwidth]{dashboard_tasks.png}
\caption{Recent Tasks and System Logs in the Dashboard}
\label{fig:dashboard_tasks}
\end{figure}

As shown in the screenshots, the dashboard captured the system processing an increasing number of URLs (from 6 to 16 to 20), all while maintaining zero failed URLs and a 0\% error rate, demonstrating the system's reliability and fault tolerance.

\begin{figure}[h]
\centering
\includegraphics[width=\textwidth]{dashboard_logs.png}
\caption{Master and Crawler Node Logs Displayed in Real-time}
\label{fig:dashboard_logs}
\end{figure}

\subsection{Search Interface Implementation}
The search interface provides an intuitive and powerful way for users to search through the indexed content. The interface is designed to support advanced search features while maintaining ease of use.

\subsubsection{Search Interface Design}
The search interface features a clean, user-friendly design with the following components:

\begin{enumerate}
    \item \textbf{Header and Search Bar}: 
    \begin{itemize}
        \item Clear "Web Crawler Search" title at the top
        \item Large search input field for query entry
        \item Prominent blue "Search" button
    \end{itemize}
    
    \item \textbf{Search Tips Section}:
    \begin{itemize}
        \item Informative tips on advanced search syntax
        \item Examples of boolean operators: "Use AND, OR, NOT for boolean search (e.g., "python AND programming")"
        \item Phrase search guidance: "Use quotes for exact phrases (e.g., "machine learning")"
        \item Field-specific search examples: "Use field-specific search (e.g., title:github, content:"open source")"
    \end{itemize}
    
    \item \textbf{Search Results Area}:
    \begin{itemize}
        \item Clean presentation of search results
        \item Each result includes title, URL, summary, keywords, and relevance score
        \item Results are ordered by relevance
    \end{itemize}
\end{enumerate}

\begin{figure}[h]
\centering
\includegraphics[width=\textwidth]{search_interface.png}
\caption{Web Crawler Search Interface}
\label{fig:search_interface}
\end{figure}

\subsubsection{Advanced Search Capabilities}
The search interface demonstrates several advanced search features in action:

\begin{enumerate}
    \item \textbf{Boolean Search}: The interface supports complex queries using AND, OR, and NOT operators, as shown in the screenshot where "python AND github" is used to find pages that contain both terms.
    
    \item \textbf{Field-Specific Search}: Users can target specific fields with queries like:
    \begin{itemize}
        \item "title:python" to search only in page titles
        \item "content:"open source"" to find exact phrases in the content
    \end{itemize}
    
    \item \textbf{Exact Phrase Matching}: The interface supports exact phrase matching using quotes, as shown in the screenshots with searches like "content:"open source"".
\end{enumerate}

\begin{figure}[h]
\centering
\includegraphics[width=\textwidth]{boolean_search.png}
\caption{Boolean Search Example with Advanced Query Options}
\label{fig:boolean_search}
\end{figure}

\subsubsection{Search Result Presentation}
Search results are presented with rich metadata to help users quickly find relevant information:

\begin{enumerate}
    \item \textbf{Result Pages}: Each search result includes:
    \begin{itemize}
        \item Clickable title linking to the original URL
        \item Full URL displayed below the title
        \item Text summary or excerpt from the page content
        \item Extracted keywords displayed as tags (e.g., "python, package, pypi, software, find, install, index, search, project, learn")
        \item Relevance score indicating match quality (e.g., "3.97")
        \item Last updated timestamp (e.g., "2025-04-26 03:44:05")
    \end{itemize}
    
    \item \textbf{Result Organization}: 
    \begin{itemize}
        \item Results are presented in relevance order (highest scores first)
        \item Clear visual separation between results with padding and border styling
        \item Keywords are highlighted in green to stand out
    \end{itemize}
\end{enumerate}

\begin{figure}[h]
\centering
\includegraphics[width=\textwidth]{search_results.png}
\caption{Search Results with Metadata and Relevance Scores}
\label{fig:search_results}
\end{figure}

The screenshots demonstrate successful searches for various queries, including:
\begin{itemize}
    \item "python AND github" - Finding pages that mention both Python and GitHub
    \item "title:python" - Finding pages with Python in the title
    \item "software" - General search for software-related content
    \item "content:"open source"" - Finding pages containing the exact phrase "open source"
\end{itemize}

\begin{figure}[h]
\centering
\includegraphics[width=\textwidth]{field_search.png}
\caption{Field-Specific Search Example}
\label{fig:field_search}
\end{figure}

The search results show rich content including GitHub Actions documentation, Wikipedia pages, PyPI (Python Package Index) documentation, and Python Software Foundation information, demonstrating the effectiveness of the search engine and interface.

\section{References}
\begin{enumerate}
    \item Whoosh Documentation - \url{https://whoosh.readthedocs.io/}
    \item NLTK Documentation - \url{https://www.nltk.org/}
    \item AWS S3 Documentation - \url{https://docs.aws.amazon.com/s3/}
    \item Flask Documentation - \url{https://flask.palletsprojects.com/}
    \item MPI for Python Documentation - \url{https://mpi4py.readthedocs.io/}
    \item Bootstrap Documentation - \url{https://getbootstrap.com/docs/}
    \item Chart.js Documentation - \url{https://www.chartjs.org/docs/}
\end{enumerate}

\end{document} 